% !TEX root = ../main_zh.tex
\chapter{增长的微观基础：新古典增长模型}
\label{ch:neoclassical}

第~\ref{ch:solow} 章的索罗模型带我们走了很远：仅凭一条资本积累方程，它就解释了为什么穷国比富国增长得快、为什么经济会收敛到稳态、以及为什么长期增长归根结底要靠技术进步而非单靠节俭。但这份解释力是有代价的。储蓄率 $s$ 完全是从模型之外硬塞进来的，是一个靠假设固定下来的常数。一个无论如何都把固定比例的收入存起来的经济，本质上就是一个\emph{计划}经济：没有人被问及存这么多是否符合自身利益，模型也对由此得到的资本存量究竟是偏多、偏少还是恰到好处保持沉默。

新古典增长模型补上了索罗所缺的那一块——\emph{微观基础}。我们让经济中住进做最优化决策的主体：家庭选择消费与储蓄以最大化效用，厂商选择投入以最大化利润，并要求各市场出清。于是储蓄率不再是我们外生强加的参数，而是家庭在市场给定的回报下、在今天的消费与明天的消费之间权衡所\emph{内生得出的结果}。值得玩味的是，当我们把这套最优化推演到底，得到的长期预言竟与索罗的几乎一模一样：稳定的稳态、收敛、以及由技术驱动的增长。不同之处在于，如今这些预言扎根于个体选择，因此我们可以提出索罗根本无从提起的福利问题——竞争性结果是否有效率？计划者能否做得更好？——并给出回答。

我们分阶段搭建这套工具。先从一个静态一般均衡经济入手，在最简单的设定里把家庭最优化、厂商最优化与市场出清的逻辑钉死。然后引入时间：一个两期禀赋经济带出欧拉方程与消费平滑；一个世代交叠经济把这套两期逻辑变成一套真正的资本积累理论，配上运动方程与稳态；最后，一个无穷期的吃蛋糕问题及其生产版本——社会计划者问题——补全全局，让我们得以陈述把竞争与效率绑在一起的两条福利定理。

\section{静态一般均衡模型}

我们从完全没有资本积累的情形起步。在单一一期之内资本存量无法改变——投资需要时间才能转化为生产能力——所以我们把 $K$ 固定住，研究在这一期之内代表性家庭与代表性厂商如何在商品与劳动两个市场上相互作用。这就把问题剥到只剩它的一般均衡内核：每一侧都是做最优化的主体，价格不断调整直到供给等于需求。

\asm{静态新古典环境}{
\begin{itemize}
  \item 各主体彼此相同，因此我们可以只研究单一的\emph{代表性}家庭与单一的\emph{代表性}厂商。
  \item 家庭在给定价格下最大化效用；厂商在给定价格下最大化利润。
  \item 市场出清：在每个市场上，需求等于供给。
  \item 资本存量 $K$ 在本期内是外生的（它是一个状态变量，不是选择变量）。
\end{itemize}
}

一开始就值得指出，宏观经济学在侧重点上将与读者在微观经济学里见过的有所不同。家庭的两类选择支配着此后的一切。\emph{劳动与闲暇}之间的选择——工作多少——本质上是\emph{静态}的：它在单一一期之内就解决了。\emph{消费与储蓄}之间的选择——为未来留下多少——本质上是\emph{动态}的，而正是这一选择，一期接一期地堆积起资本存量 $K$。静态模型把第一类选择单独拎出来；本章其余部分则展开第二类。

\subsection{家庭：消费与闲暇}

代表性家庭从消费 $C$ 与闲暇 $l$ 中获得效用，
\[
U = U(C, l),
\]
我们假定通常的良态条件：两种物品都是值得拥有的，$U_C > 0$ 且 $U_l > 0$；边际替代率递减（无差异曲线是凸的）；并且消费与闲暇都是正常品。

家庭被赋予固定的时间禀赋 $H$，把它在闲暇 $l$ 与工作 $H - l$ 之间分配。我们把消费品的价格标准化为一，并令 $w$ 表示实际工资，于是 $w$ 同时也是闲暇的价格——每一小时不工作，家庭就要付出 $w$ 单位被放弃的消费作为代价。除工资收入之外，家庭还获得厂商分配的利润 $\pi$，并缴纳一笔总额税 $T$。其预算约束为
\[
C = w(H - l) + \pi - T.
\]
把闲暇的工资价值移到左边，约束就呈现出它最有启发性的形式，
\[
C + w\,l = wH + \pi - T,
\]
这读作一个\emph{完全收入}（full income）预算：家庭的全部资源，即右边，等于其全部时间禀赋的价值加上利润减税之后的余额，而它把这些资源花在消费上、以及花在以 $w$ 计价的闲暇上。右边是一个家庭视为既定的数。

在此约束下最大化 $U(C,l)$，标准的相切条件是闲暇与消费之间的边际替代率等于闲暇的相对价格，
\[
\mathrm{MRS} = w \quad\Longrightarrow\quad C^*,\; l^*.
\]
于是家庭的最优劳动供给为 $N^s = H - l^*$。

\subsection{厂商：生产与劳动需求}

代表性厂商用资本与劳动生产产出，
\[
Y = z\,F(K, N^d),
\]
其中 $z$ 是全要素生产率，$N^d$ 是雇佣的劳动。我们假定 $F$ 呈现规模报酬不变（CRS），正因如此我们才可以把整个生产侧当作单一的价格接受厂商来处理：在 CRS 下单个厂商的规模无从确定，而总量层面表现得就像一个代表性生产者。技术具有为正但递减的边际产出，
\[
F_K > 0, \quad F_N > 0, \qquad F_{KK} < 0, \quad F_{NN} < 0.
\]

\rmkb[对课堂讲义的一处更正]{讲义把二阶条件写成了 $F_{KK}>0$ 与 $F_{NN}>0$。这是笔误：边际产出递减——正是这一特征让模型表现良好、并在下文的增长版本里给出收敛——要求\emph{二}阶导数为\emph{负}，即 $F_{KK}<0$ 与 $F_{NN}<0$。一阶导数为正；二阶导数为负。}

在本期之内，$K$ 是一个固定的状态变量，所以厂商只选择雇佣多少劳动。它的问题是
\[
\max_{N^d}\; \set{ z\,F(K, N^d) - w\,N^d }.
\]
一阶条件令劳动的边际产出等于实际工资，
\[
\mathrm{MPN} = z\,F_N(K, N^d) = w \quad\Longrightarrow\quad N^{d*}.
\]

\rmk{注意这两个条件之间的分工：$\mathrm{MRS}=w$ 是\emph{家庭}的最优化条件，$\mathrm{MPN}=w$ 是\emph{厂商}的。课堂上的速记“$\mathrm{MPN}=\mathrm{MRS}=w$”把两者压成了一行，但它们分属不同的主体；只有\emph{在均衡中}、当唯一的工资 $w$ 调整到让两侧都满意时，它们才同时成立。}

\subsection{均衡与瓦尔拉斯定律}

一个均衡就是一个让两个市场都出清的工资 $w$。在劳动市场上，家庭的供给等于厂商的需求，
\[
H - l^* = N^{d*}.
\]
在商品市场上，家庭消费的恰好就是厂商生产的，
\[
C^* = Y^*.
\]

这两个出清条件并不相互独立。由家庭的预算约束可知，它所需求之物的价值始终等于它所供给之物的价值；把各市场加总，超额需求必然彼此抵消。这就是\emph{瓦尔拉斯定律}（Walras' law）：在一个有两个市场的经济里，只要一个出清，另一个也必然出清。

\state{瓦尔拉斯定律}{
在一个有 $n$ 个市场的一般均衡中，总超额需求的价值恒等于零。因此，只要 $n-1$ 个市场出清，最后一个市场就自动出清。在我们这个两市场经济里，\textbf{一旦劳动市场出清，商品市场也随之出清}——我们只需找到那个使劳动供给等于劳动需求的唯一工资 $w$ 即可。}

可以证明，存在一个\emph{唯一}的工资 $w$ 使得 $H - l^* = N^{d*}$。在该工资下，厂商的利润 $\pi$ 被确定，家庭的完全收入被确定，整个配置 $(C^*, l^*, N^{d*}, Y^*)$ 也随之落定。

\subsection{工资变动的收入效应与替代效应}

均衡对工资变动如何反应？由于工资既是家庭每工作一小时的收入、又是闲暇的价格，$w$ 的变动会引发两股相互角力的力量。

设想 $w$ 上升。\emph{收入效应}：家庭如今更富了，又因为消费与闲暇都是正常品，它两者都想要更多——消费上升，闲暇也上升（工作减少）。\emph{替代效应}：闲暇相对于消费变贵了，于是家庭从闲暇转向消费——消费上升，闲暇下降（工作增加）。

两种效应把消费推向同一方向，所以工资上升\emph{毫无疑义地抬高消费}。但它们把闲暇——从而把劳动供给——推向相反方向，所以对工作小时数的净效应一般是不确定的。劳动供给随工资上升还是下降，取决于哪种效应占上风。

有一个干净的特例让这种抵消恰好成立。把时间禀赋标准化为一（$H=1$，于是当 $\pi=T=0$ 时 $C = w(1-l)$），并取对数偏好，
\[
\max_{C,\,l}\; \log C + \beta \log l
\qquad \text{s.t.} \qquad C + w\,l = w.
\]

\ex[对数效用：收入效应与替代效应相互抵消]{
求解家庭的问题，并证明劳动供给与工资无关。}
\sol{
把预算约束 $C = w(1-l)$ 代入目标函数，使之成为仅关于 $l$ 的函数：
\[
U(l) = \log\!\big(w(1-l)\big) + \beta \log l = \log w + \log(1-l) + \beta \log l.
\]
一阶条件为
\[
\dd{U}{l} = \frac{-1}{1-l} + \frac{\beta}{l} = 0
\quad\Longrightarrow\quad
\beta(1-l) = l
\quad\Longrightarrow\quad
l^* = \frac{\beta}{1+\beta}.
\]
于是劳动供给为
\[
1 - l^* = \frac{1}{1+\beta},
\]
这是一个常数：工资 $w$ 已彻底退出。在对数效用下，闲暇上的收入效应与替代效应恰好相互抵消，所以工作小时数与工资无关。工资上升的全部作用都流向了消费：由预算约束，$C^* = w(1-l^*) = w/(1+\beta)$ 随 $w$ 一比一地上升。}

这一刀刃般的结论并非代数上的巧合；对数效用恰恰是替代弹性与收入弹性调得使两种效应相互抵消的那个特例。这正是对数效用在本章里反复出现的原因：它给出最简单的闭式解，把我们关心的动态机制从恼人的、作用于劳动的财富效应中分离出来。

\section{两期禀赋经济}

现在我们以最简单的形式引入时间。暂且把生产和厂商剥离，考虑一个恰好活两期、且每一期都获得一笔外生收入的家庭——一份天上掉下来的禀赋（windfall），像吗哪一样。这里没有什么可生产，也没有劳动可供给；唯一的决策是\emph{何时}消费。仅凭这副最朴素的设定，就足以引出动态宏观经济学的核心对象——欧拉方程。

家庭最大化终身效用
\[
\max_{C_1,\,C_2}\; U(C_1, C_2) = \log C_1 + \beta \log C_2,
\]
其中 $\beta \in (0,1)$ 是\emph{贴现因子}（discount factor）——一个接近 $0.95$ 的数，刻画了人的不耐：明天的一单位效用不如今天的一单位值钱。

\defn{贴现因子}{
\emph{贴现因子} $\beta \in (0,1)$ 是家庭对下一期效用相对于本期效用所赋予的权重。$\beta$ 越大表示越有耐心。与之对应的\emph{贴现率} $\rho$ 满足 $\beta = 1/(1+\rho)$。}

收入在第 $1$ 期到来为 $Y_1$，在第 $2$ 期到来为 $Y_2$，二者皆为外生。家庭只能通过以市场利率 $r$ 储蓄一笔金额 $S$ 来在两期之间转移资源。它逐期的预算约束为
\[
\ba
C_1 + S &= Y_1, \\
C_2 &= Y_2 + (1+r)\,S.
\ea
\]
第一期的储蓄在第二期赚得毛回报 $1+r$。和静态模型一样，从家庭的视角看我们把 $r$、$Y_1$、$Y_2$ 都当作既定：决定利率的宏观力量超出了这单个家庭的控制范围。

储蓄 $S$ 是联结两期的桥梁。把它消去——从第二个约束解出 $S$ 代入第一个——便把两条单期预算约束合并成单一的\emph{跨期预算约束}，
\[
C_1 + \frac{C_2}{1+r} = Y_1 + \frac{Y_2}{1+r}.
\]
其含义正是人们所期望的：终身消费的现值等于终身收入的现值。

\subsection{欧拉方程}

在最优处，家庭无法通过把边际一单位消费在两期之间挪动来提高效用。相切条件是跨期边际替代率等于毛利率——即以明天的消费衡量的、今天消费的相对价格：
\[
\mathrm{MRS} = -\dd{C_2}{C_1}\bigg|_{U} = \frac{\mathrm{MU}_1}{\mathrm{MU}_2} = 1 + r.
\]
在 $U = \log C_1 + \beta \log C_2$ 下我们有 $\mathrm{MU}_1 = 1/C_1$ 与 $\mathrm{MU}_2 = \beta/C_2$，于是
\[
\frac{1/C_1}{\beta/C_2} = \frac{C_2}{\beta C_1} = 1 + r,
\]
整理后即得\emph{欧拉方程}（Euler equation）。

\thm{消费的欧拉方程（两期，对数效用）}{
最优消费路径满足
\[
\frac{C_2^*}{C_1^*} = \beta(1+r),
\qquad\text{等价地}\qquad
C_2^* = \beta(1+r)\,C_1^*.
\]}

\pf{
由跨期预算约束 $C_1 + C_2/(1+r) = W$，其中 $W := Y_1 + Y_2/(1+r)$ 是终身财富，拉格朗日函数为
\[
\mathcal{L} = \log C_1 + \beta \log C_2 + \lambda\!\pr{W - C_1 - \frac{C_2}{1+r}}.
\]
一阶条件为
\[
\frac{1}{C_1} = \lambda, \qquad \frac{\beta}{C_2} = \frac{\lambda}{1+r}.
\]
第一式除以第二式消去 $\lambda$，得 $\dfrac{C_2}{\beta C_1} = 1+r$，即 $C_2^* = \beta(1+r)\,C_1^*$。
}

不构造跨期预算约束，也能得到同一个条件：直接把约束代入并对 $S$ 求最大，
\[
U(S) = \log\!\big(Y_1 - S\big) + \beta \log\!\big(Y_2 + (1+r)S\big),
\qquad
\frac{-1}{Y_1 - S} + \frac{\beta(1+r)}{Y_2 + (1+r)S} = 0,
\]
回代 $C_1 = Y_1 - S$ 与 $C_2 = Y_2 + (1+r)S$ 后，依然得到 $C_2 = \beta(1+r)C_1$。

\subsection{消费平滑}

欧拉方程是理解家庭如何把消费铺展到时间之上的钥匙。假定有那么一刻既无不耐也无利息，即 $\beta = 1$ 且 $r = 0$。那么欧拉方程读作 $C_2^* = C_1^*$：家庭选择在两期消费完全相同的量，无论其收入的时间分布如何。这就是\emph{消费平滑}（consumption smoothing）。

\state{为什么凹性意味着平滑}{
只要效用是严格凹的（边际效用为正且递减，$U'>0$，$U''<0$），家庭就偏好一条平滑的消费路径，而非一条现值相同却起伏的路径。从高消费的一期取走一单位消费、挪到低消费的一期，会提高效用，因为失去的边际效用小于得到的边际效用。平滑正是边际效用递减在行为上的内涵。}

一旦把不耐与利息重新放回，路径会倾斜，但不会断裂。欧拉方程 $C_2^* = \beta(1+r)C_1^*$ 说家庭是\emph{围绕}因子 $\beta(1+r)$ 来平滑的。直觉正如人们所料：把一单位储蓄留到下一期会赚得利息，资源因而增长 $(1+r)$ 倍，这把消费向未来倾斜；但下一期的效用要被 $\beta$ 折现，这又把它拉回当下。当 $\beta(1+r) = 1$ 时这两股力量平衡，消费持平；当 $\beta(1+r) > 1$ 时回报盖过不耐，消费随时间上升；当 $\beta(1+r) < 1$ 时不耐取胜，消费随时间下降。

\subsection{值函数与影子价格}

一旦参数 $(Y_1, Y_2, r)$ 被固定下来，这场最优化就钉死了唯一的解 $(C_1^*, C_2^*, S^*)$，从而也钉死了唯一的最大化效用。把这个最大化效用看作那些参数的函数，它就是\emph{值函数}（value function），
\[
U^* = V(Y_1, Y_2, r).
\]
值函数记录了在给定处境下家庭所能达到的最好结果。它的导数携带着经济含义。对第一期收入求导，$\partial V / \partial Y_1$，衡量第 $1$ 期多一单位天降收入会把终身福利抬高多少；这恰恰是第 1 期资源的\emph{影子价格}（shadow price）——即最优化里的拉格朗日乘子 $\lambda$。在最优处，乘子等于每种物品的边际效用与其价格之比，
\[
\frac{\partial U / \partial x_i}{p_i} = \lambda \qquad \text{对每个选择 } x_i,
\]
这正是那个熟悉的论断：在均衡中，家庭把每一元钱的边际效用在所有边际上拉平。下文的无穷期问题里我们将再次遇见值函数及其影子价格，到那时这种递归的思考方式将变得不可或缺。

\section{世代交叠模型}

两期经济教会了我们单个家庭如何把消费配置到时间之上，但它没有生产、没有资本，所以还称不上一套\emph{增长}理论。世代交叠（OLG）模型把这一点补上：它把两期的人生一个叠一个地摞起来，使得在任何时点上，年轻一代与年老一代都并存。年轻人工作并储蓄；他们的储蓄变成下一期所继承的资本。如此一来，两期模型里那个动态的储蓄决策就成了资本积累的引擎。

\begin{figure}[h]
\centering
\begin{tikzpicture}[scale=1.0]
  % --- geometry parameters ---
  \def\bw{2.0}   % box width  (one period)
  \def\bh{0.9}   % box height (one generation level)
  \def\ng{4}     % number of overlapping generations
  % The bottom of generation g sits at level g (g=0,...,3),
  % and its "Work" box starts at x = g*\bw.  Each generation
  % occupies two consecutive period-boxes: Work then Retire.
  %
  % --- period-boundary dashed guide lines ---
  % boundaries fall at x = 1,2,3,4 periods (in box-width units).
  \pgfmathsetmacro{\xa}{1*\bw}
  \pgfmathsetmacro{\xb}{2*\bw}
  \pgfmathsetmacro{\xc}{3*\bw}
  \pgfmathsetmacro{\xd}{4*\bw}
  \pgfmathsetmacro{\toptop}{\ng*\bh+0.55}
  \foreach \xx in {\xa,\xb,\xc,\xd}{
    \draw[densely dashed] (\xx,-0.15) -- (\xx,\toptop);
  }
  % --- the staircase of generations ---
  \foreach \g in {0,1,2,3}{
    \pgfmathsetmacro{\xstart}{\g*\bw}
    \pgfmathsetmacro{\ylo}{\g*\bh}
    \pgfmathsetmacro{\yhi}{\g*\bh+\bh}
    \pgfmathsetmacro{\xmid}{\xstart+\bw}
    \pgfmathsetmacro{\xend}{\xstart+2*\bw}
    % Work box
    \draw (\xstart,\ylo) rectangle (\xmid,\yhi);
    \node at ({\xstart+0.5*\bw},{\ylo+0.5*\bh}) {Work};
    % Retire box
    \draw (\xmid,\ylo) rectangle (\xend,\yhi);
    \node at ({\xmid+0.5*\bw},{\ylo+0.5*\bh}) {Retire};
  }
  % --- Time axis (arrow) along the bottom ---
  \draw[-Latex,line width=1.1pt] (-0.2,-0.55) -- ({5*\bw+0.4},-0.55)
        node[below right] {Time};
\end{tikzpicture}

\caption{世代交叠的时间线：每一代年轻时工作并储蓄，年老时靠积累下来的储蓄度日，于是在每个时点上，一代正在工作的年轻人都与一代已退休的年老人交叠在一起。}
\label{fig:macro-5}
\end{figure}

这一结构如图~\ref{fig:macro-5} 所示。每个个体活两期。年轻时她供给一单位劳动、赚取工资，消费其中一部分、把其余储蓄起来；年老时她不再工作，靠储蓄度日（此时储蓄已化为她租给厂商的资本）。尽管任何一个人都只活两期，世代交叠的链条却无尽地延伸下去，所以这个模型实际上是一个\emph{无穷期}经济——只是我们一次只截取一代来看罢了。

\subsection{家庭}

一个在 $t$ 期出生的个体，年轻时供给一单位劳动，故其工资收入为 $w_t$。她的预算约束为
\[
\ba
C_t + S_t &= w_t, \\
C_{t+1} &= (1+r_{t+1})\,S_t,
\ea
\]
其中 $S_t$ 是储蓄，$r_{t+1}$ 是这笔储蓄所赚得的净回报。关键的会计联结是：一代人的储蓄以资本的形式持有，所以年轻人的储蓄变成下一期的资本存量，
\[
S_t = K_{t+1}.
\]
正是这一个方程，把一连串两期问题变成了一套积累理论：明天社会动用的资本，恰恰就是今天年轻人选择留出来的那部分。

在对数偏好下，
\[
U(C_t, C_{t+1}) = \ln C_t + \beta \ln C_{t+1},
\]
储蓄决策就是我们刚刚解过的那个问题。欧拉方程 $C_{t+1}/C_t = \beta(1+r_{t+1})$ 连同预算约束，给出一条明确的储蓄法则。

\ex[OLG 模型中的储蓄]{
在对数效用与上述预算约束下，证明年轻人储蓄的是工资的一个固定比例，从而
\[
K_{t+1} = \frac{\beta}{1+\beta}\, w_t.
\]}
\sol{
把约束代入目标函数。在 $C_t = w_t - S_t$ 与 $C_{t+1} = (1+r_{t+1})S_t$ 下，
\[
U(S_t) = \ln(w_t - S_t) + \beta \ln\!\big((1+r_{t+1})S_t\big).
\]
关于 $S_t$ 的一阶条件为
\[
\frac{-1}{w_t - S_t} + \frac{\beta}{S_t} = 0
\quad\Longrightarrow\quad
\beta(w_t - S_t) = S_t
\quad\Longrightarrow\quad
S_t = \frac{\beta}{1+\beta}\, w_t.
\]
毛回报 $1+r_{t+1}$ 被消掉了——这又是对数效用的一个后果——所以储蓄是工资的固定比例 $\beta/(1+\beta)$，与利率无关。由于 $S_t = K_{t+1}$，
\[
K_{t+1} = \frac{\beta}{1+\beta}\, w_t.
\]}

我们保留 $K$ 而不把它代换掉，因为它是联结市场两侧的那个变量——家庭通过储蓄供给它，厂商为生产而需求它——它同时也是联结相邻两期的那个变量。

\subsection{厂商}

厂商每期租用资本、雇佣劳动以最大化利润，
\[
\max_{K,\,L}\; F(K, L) - w\,L - r\,K,
\]
采用柯布--道格拉斯的、劳动增进型技术
\[
F(K, L) = K^{\alpha}\,(zL)^{1-\alpha}, \qquad 0 < \alpha < 1.
\]
一阶条件令每种要素的价格等于其边际产出，
\[
\ba
w_t &= (1-\alpha)\,z_t^{\,1-\alpha}\, K_t^{\alpha}\, L_t^{-\alpha}, \\
r_t &= \alpha\,z_t^{\,1-\alpha}\, K_t^{\alpha-1}\, L_t^{\,1-\alpha}.
\ea
\]

\subsection{市场均衡与运动方程}

在均衡中，每种要素的价格在供给侧与需求侧是相同的，所以我们可以把厂商的工资代入家庭的储蓄法则。结果是一个关于资本的差分方程，
\[
K_{t+1} = \frac{\beta}{1+\beta}\, w_t = \frac{\beta}{1+\beta}\,(1-\alpha)\,z_t^{\,1-\alpha}\, L_t^{-\alpha}\, K_t^{\alpha}.
\]
由于问题里处处都没有出现人口变量、且我们已把每个年轻人的劳动供给标准化为一单位，我们令 $L_t = 1$，并把每个量都读作人均量。暂且把生产率固定在常数水平 $z_t = z$，运动方程就变成
\[
K_{t+1} = \frac{\beta(1-\alpha)}{1+\beta}\, z^{\,1-\alpha}\, K_t^{\alpha}.
\]

\begin{figure}[h]
\centering
\begin{tikzpicture}[scale=1.0]
  % ---- parameters of the law of motion K_{t+1} = B * K_t^alpha ----
  \def\al{0.5}             % concavity exponent (alpha < 1)
  \def\Kss{5.2}            % steady-state capital (curve meets 45-degree line)
  \pgfmathsetmacro{\B}{\Kss/pow(\Kss,\al)}   % B so that B*Kss^alpha = Kss
  \def\xmax{8.0}
  \def\ymax{7.2}
  % ---- axes ----
  \draw[-Latex] (0,0) -- (\xmax,0) node[right] {$K_t$};
  \draw[-Latex] (0,0) -- (0,\ymax) node[above] {$K_{t+1}$};
  % ---- 45-degree line K_{t+1} = K_t (dotted, as in the original) ----
  \draw[densely dotted,line width=1.0pt] (0,0) -- (6.5,6.5);
  \node[above left=1pt and -1pt] at (6.2,6.2) {$45^{\circ}$};
  % ---- law of motion curve K_{t+1} = B*K_t^alpha (concave, through origin) ----
  \draw[line width=1.1pt,smooth,samples=160,domain=0:6.9]
        plot (\x,{\B*pow(\x,\al)});
  \node[right] at (5.6,4.55) {$\dfrac{\beta}{1-\beta}\,(1-\alpha)\,z\,K_t^{\alpha}$};
  % ---- steady state where curve meets the 45-degree line ----
  \fill (\Kss,\Kss) circle (1.8pt);
  \draw[densely dashed] (\Kss,\Kss) -- (\Kss,0) node[below] {$K_{SS}$};
  \draw[densely dashed] (\Kss,\Kss) -- (0,\Kss) node[left] {$K_{SS}$};
\end{tikzpicture}

\caption{资本的运动方程 $K_{t+1}=f(K_t)$。由于 $\alpha<1$，曲线是凹的，且与 $45^\circ$ 线恰好相交一次于稳态 $K^{*}$；从任何为正的初始点出发，经济都收敛到该点。}
\label{fig:macro-6}
\end{figure}

由于指数 $\alpha$ 小于一，右边是 $K_t$ 的一个凹的、递增的函数，并且与 $45^\circ$ 线恰好相交一次（图~\ref{fig:macro-6}）。这个交点就是稳态，在那里资本自我复制，$K_{t+1} = K_t = K_{\mathrm{ss}}$。

\ex[稳态资本]{
从运动方程解出 $K_{\mathrm{ss}}$。}
\sol{
在稳态处 $K_{\mathrm{ss}} = \dfrac{\beta(1-\alpha)}{1+\beta}\, z^{\,1-\alpha}\, K_{\mathrm{ss}}^{\alpha}$。两边同除以 $K_{\mathrm{ss}}^{\alpha}$，
\[
K_{\mathrm{ss}}^{\,1-\alpha} = \frac{\beta(1-\alpha)}{1+\beta}\, z^{\,1-\alpha}
\quad\Longrightarrow\quad
K_{\mathrm{ss}} = z\pr{\frac{\beta(1-\alpha)}{1+\beta}}^{\frac{1}{1-\alpha}}.
\]}

通过一期接一期地施用运动方程，我们实际上已把两期的 OLG 结构延展到了无穷多期。又由于没有任何人口项出现，每个变量都带有人均的含义。运动方程还重现了索罗的那些教训：一个起步资本很少的国家位于图~\ref{fig:macro-6} 中靠左的位置，那里曲线很陡，所以它的资本——以及产出——增长得很快；相似的经济收敛到相同的稳态。其深层原因，再一次，是 $\alpha < 1$ 所编码的资本边际产出递减。

\subsection{稳定稳态与不稳定稳态}

值得就差分方程的一个一般要点停一停，因为并非每个不动点都是归宿。考虑任意一阶差分方程
\[
x_{t+1} = f(x_t), \qquad t = 1, 2, \dots
\]
这就是那个典范的\emph{运动方程}，或曰马尔可夫链结构：明天的取值只取决于今天的取值，而不取决于全部历史。稳态就是一个不动点，
\[
x_{\mathrm{ss}} = f(x_{\mathrm{ss}}),
\]
即 $f$ 的图像与 $45^\circ$ 线相交之处。但解出 $x = f(x)$ 只告诉我们稳态\emph{存在}；它并不告诉我们经济是否真的会朝它移动。

\begin{figure}[h]
\centering
\begin{tikzpicture}[scale=1.0]
  % axes
  \draw[-Latex] (0,0) -- (5.6,0) node[right] {$x_t$};
  \draw[-Latex] (0,0) -- (0,5.4) node[above] {$x_{t+1}$};
  % 45-degree line x_{t+1}=x_t (dotted)
  \draw[densely dotted] (0,0) -- (5.0,5.0);
  \node[above] at (5.0,5.0) {$45^{\circ}$};
  % law of motion f: concave, crosses the 45-line from above at x*=3
  % f(x) = 3*sqrt(x/3) = sqrt(3)*sqrt(x); f(3)=3, f'(3)=1/2<1 (stable)
  \draw[line width=1.0pt,smooth,samples=160,domain=0.02:5.1]
        plot (\x,{sqrt(3)*sqrt(\x)});
  \node[right] at (5.1,{sqrt(3)*sqrt(5.1)}) {$f(x_t)$};
  % steady state at the crossing
  \fill (3,3) circle (1.7pt);
  \draw[densely dashed] (3,3) -- (3,0) node[below] {$x^{*}$};
  \draw[densely dashed] (3,3) -- (0,3) node[left] {$x^{*}$};
\end{tikzpicture}

\caption{一个\emph{稳定}的稳态：运动方程穿过 $45^\circ$ 线时斜率的绝对值小于一，所以在附近起步的轨迹会被拉回到交点。}
\label{fig:macro-7}
\end{figure}

\begin{figure}[h]
\centering
\begin{tikzpicture}[scale=1.0]
  % axes
  \draw[-Latex] (0,0) -- (5.6,0) node[right] {$x_t$};
  \draw[-Latex] (0,0) -- (0,5.0) node[above] {$x_{t+1}$};
  % 45-degree line x_{t+1}=x_t
  \draw[densely dashed] (0,0) -- (4.6,4.6) node[above right] {$45^{\circ}$};
  % law of motion f: convex, steeper than 45 line at the crossing
  % f(x)=x^2/x* with x*=2.2 -> crosses 45-line at (x*,x*) with slope 2>1
  \def\xs{2.2}
  \draw[line width=1.1pt,smooth,samples=120,domain=0:3.18,variable=\x]
        plot ({\x},{\x*\x/\xs}) node[right] {$f(x_t)$};
  % steady state (unstable): crossing of f with the 45-degree line
  \fill (\xs,\xs) circle (1.7pt);
  \draw[densely dashed] (\xs,\xs) -- (\xs,0) node[below] {$x_{SS}$};
  \draw[densely dashed] (\xs,\xs) -- (0,\xs) node[left] {$x_{SS}$};
\end{tikzpicture}

\caption{一个\emph{不稳定}的稳态：运动方程穿过 $45^\circ$ 线时斜率大于一，所以交点排斥而非吸引附近的轨迹。}
\label{fig:macro-8}
\end{figure}

稳定性由 $f$ 在交点处的斜率决定。若曲线从上方穿过 $45^\circ$ 线、斜率绝对值小于一（图~\ref{fig:macro-7}），轨迹会被拉回交点，稳态便是\emph{稳定}的——它是一个真正的吸引子。若曲线穿过时斜率大于一（图~\ref{fig:macro-8}），交点会\emph{排斥}附近的路径，稳态便是\emph{不稳定}的：它解出了 $x = f(x)$，但经济永远不会在那里安顿下来。我们这个凹的资本运动方程在交点处的斜率为 $\alpha\, K_{\mathrm{ss}}^{\alpha-1}\cdot(\text{常数}) < 1$，所以 $K_{\mathrm{ss}}$ 是稳定的——属于图~\ref{fig:macro-7} 那个好的情形。

\subsection{生产率增长}

至此我们把生产率冻结在了 $z_t = z$。但 $z$ 衡量的是劳动的生产效率——技术、组织、知识技能——它并不停滞不前。改设生产率以恒定速率 $g$ 增长，
\[
z_{t+1} = (1+g)\, z_t.
\]

\begin{figure}[h]
\centering
\begin{tikzpicture}[scale=1.0]
  % --- parameters: K_{t+1} = c * K_t^alpha, concave & increasing ---
  % reach at K=2 (=c*2^al): \cone->~1.13, \ctwo->~1.40, \cthree->~1.70
  \def\al{0.42}
  \def\cone{0.845}  % lowest curve  (small z)
  \def\ctwo{1.046}  % middle curve  (larger z)
  \def\cthree{1.27} % highest curve (largest z)
  % --- drawing scale: 1 unit of K -> 3.6cm ---
  \def\S{3.6}
  % --- axes ---
  \draw[-Latex] (0,0) -- ({2.18*\S},0) node[right] {$K_{t}$};
  \draw[-Latex] (0,0) -- (0,{2.18*\S}) node[above] {$K_{t+1}$};
  % axis variable captions (textbook English)
  \node[below=16pt] at ({1.09*\S},0) {Capital today};
  \node[rotate=90,above=38pt] at (0,{1.09*\S}) {Capital tomorrow};
  % --- tick marks and labels ---
  \foreach \t in {0.5,1,1.5,2}{
    \draw ({\t*\S},0) -- ({\t*\S},-0.07) node[below] {\small $\t$};
    \draw (0,{\t*\S}) -- (-0.07,{\t*\S}) node[left] {\small $\t$};
  }
  \node[below left] at (0,0) {\small $0$};
  % --- 45-degree dotted line ---
  \draw[densely dotted,thick] (0,0) -- ({2.05*\S},{2.05*\S});
  % --- three concave law-of-motion curves (black, house style) ---
  \draw[line width=1pt,smooth,samples=160,domain=0.001:2]
        plot ({\x*\S},{\cone*pow(\x,\al)*\S});
  \draw[line width=1pt,smooth,samples=160,domain=0.001:2]
        plot ({\x*\S},{\ctwo*pow(\x,\al)*\S});
  \draw[line width=1pt,smooth,samples=160,domain=0.001:2]
        plot ({\x*\S},{\cthree*pow(\x,\al)*\S});
  % --- curve labels ---
  \node[right] at ({2*\S},{\cone*pow(2,\al)*\S}) {\small $z_0$};
  \node[right] at ({2*\S},{\ctwo*pow(2,\al)*\S}) {\small $z_1$};
  \node[right] at ({2*\S},{\cthree*pow(2,\al)*\S}) {\small $z_2$};
  % --- blue accent arrows: steady state migrates outward as z grows ---
  % lower arrow near middle-curve crossing of 45-line (~K=1.05)
  \draw[-Latex,line width=1.6pt,blue!75!black] ({1.02*\S},{0.94*\S}) -- ({1.16*\S},{1.12*\S});
  % upper arrow near top-curve crossing of 45-line (~K=1.5)
  \draw[-Latex,line width=1.6pt,blue!75!black] ({1.40*\S},{1.27*\S}) -- ({1.56*\S},{1.47*\S});
\end{tikzpicture}

\caption{生产率增长使运动方程逐期上移：随着 $z_t$ 上升，曲线 $K_{t+1}=f(K_t)$ 向外移动，拖着那个移动的目标 $K_{\mathrm{ss}}(z_t)$ 稳步右移，并把资本存量沿一条增长路径拉着走。}
\label{fig:macro-9}
\end{figure}

运动方程仍读作
\[
K_{t+1} = \frac{\beta}{1+\beta}\, w_t = \frac{\beta(1-\alpha)}{1+\beta}\, z_t^{\,1-\alpha}\, L_t^{-\alpha}\, K_t^{\alpha},
\]
若令 $L_t = 1$ 并暂时把 $z_t$ 当作固定，与此前同样的代数给出一个目标水平
\[
K_{\mathrm{ss}} = z_t\pr{\frac{\beta(1-\alpha)}{1+\beta}}^{\frac{1}{1-\alpha}}.
\]
但如今 $z_t$ 自身在增长，所以严格地说 $K_{t+1} \ne f(K_t)$ 并不对应一个不随时间变化的 $f$，上面那个“$K_{\mathrm{ss}}$”也就不是一个真正的稳态。它毋宁是一个\emph{移动的目标}：一条随 $z_t$ 上升而向上漂移的增长路径（图~\ref{fig:macro-9}）。若资本存量偏离了这条路径，边际报酬递减那股力量仍会把它拉回路径之上。于是在长期，资本以与生产率相同的速率增长，
\[
K_{t+1} = (1+g)\, K_t,
\]
因为 $K_{\mathrm{ss}}$ 与 $z_t$ 成比例。

把这种平衡增长代入生产函数，取 $L_t = 1$，
\[
Y_{t+1} = K_{t+1}^{\alpha}\,(z_{t+1} L_{t+1})^{1-\alpha}
        = \big[(1+g)K_t\big]^{\alpha}\big[(1+g)z_t L_{t+1}\big]^{1-\alpha}
        = (1+g)\, Y_t,
\]
所以产出也以生产率的增速 $g$ 增长。长期增长由技术驱动，与索罗一模一样。

$Y$ 与 $K$ 同速增长，对比值与回报率有着鲜明的含义。比值 $Y/K$ 趋于一个常数，因为两者都以 $g$ 增长。资本的回报率于是也是恒定的，
\[
r_t = \alpha\, z_t^{\,1-\alpha}\, K_t^{\alpha-1}\, L_t^{\,1-\alpha} = \frac{\alpha\, Y_t}{K_t},
\]
是那个（恒定的）产出--资本比的固定倍数。相形之下，工资随生产率增长，
\[
w_t = \frac{(1-\alpha)\,Y_t}{L_t},
\]
随人均产出的上升以速率 $g$ 上升。

\state{战后增长的卡尔多典型化事实}{
新古典模型重现了二战以来发达经济体的若干宽泛规律：
\begin{enumerate}
  \item 产出--资本比 $Y/K$ 大致恒定；
  \item 资本的实际回报率 $r$ 大致恒定；
  \item 资本的收入份额 $\alpha$ 大致恒定；
  \item 人均产出 $Y/L$ 与人均资本 $K/L$ 以共同的速率 $g$ 增长。
\end{enumerate}
索罗模型与新古典模型都与这些事实相符；不同的是，新古典模型是从最优化中赢得它们，而非靠假设给定它们。}

\section{吃蛋糕问题}

现在我们从两期人生转向一个具有无穷视野的单一决策者。\emph{吃蛋糕问题}（cake-eating problem）是最干净的无穷期最优化：一个消费者拥有一笔资源的存量——一块蛋糕——必须决定以多快的速度把它吃掉。它没有生产、没有价格，只有一笔随消费而缩小的存量，它教给我们的，是本章所有动态模型背后那套递归结构。

消费者最大化经过贴现的终身效用
\[
\max_{\set{C_t}_{t=0}^{\infty}}\; \sum_{t=0}^{\infty} \beta^t\, U(C_t)
\]
受制于资源的转移方程
\[
C_t + H_{t+1} = (1-\delta)\, H_t, \qquad \forall t,
\]
其中 $H_t$ 是第 $t$ 期初的蛋糕存量，$\delta$ 是折旧率——蛋糕每期会变质一点。这里的术语极为基本，并在动态宏观经济学中处处复现。

\defn{状态变量与控制变量}{
\emph{状态}（state）变量是决策者在一期之初所继承、且在本期之内无法改变的东西；这里它是存量 $H_t$。\emph{控制}（control）变量是她在本期之内所选择的东西；这里它是 $C_t$（等价地是 $H_{t+1}$，因为二者被资源约束联结起来）。每个变量都先后扮演这两种角色：今天的控制 $H_{t+1}$ 成为明天的状态。}

\subsection{求解该问题}

给每一期的约束附上一个乘子 $\lambda_t$，构造拉格朗日函数
\[
\mathcal{L} = \sum_{t=0}^{\infty} \beta^t\, U(C_t) + \sum_{t=0}^{\infty} \lambda_t\!\pr{(1-\delta)H_t - C_t - H_{t+1}}.
\]
关于 $C_t$、$C_{t+1}$ 以及那个起桥梁作用的存量 $H_{t+1}$ 的一阶条件为
\[
\ba
C_t:&\quad \beta^t\, U'(C_t) = \lambda_t, \\
C_{t+1}:&\quad \beta^{t+1}\, U'(C_{t+1}) = \lambda_{t+1}, \\
H_{t+1}:&\quad \lambda_{t+1}(1-\delta) = \lambda_t.
\ea
\]
存量 $H_{t+1}$ 既出现在第 $t$ 期的约束里（带负号），又出现在第 $(t+1)$ 期的约束里（乘以 $1-\delta$），这正是为什么它的一阶条件把两个乘子拴在一起。消去乘子便得到欧拉方程，
\[
\frac{U'(C_t)}{\beta\, U'(C_{t+1})} = 1 - \delta.
\]
特化到对数效用 $U(C_t) = \log C_t$，于是 $U'(C) = 1/C$，
\[
\frac{C_{t+1}}{C_t} = \beta(1-\delta).
\]

\defn{政策函数}{
\emph{政策函数}（policy function）把当前状态映射为最优控制——它勾勒出经过最优化的决策实际所沿行的那条路径。这里它把消费表示为所继承存量的函数。}

欧拉方程 $C_{t+1} = \beta(1-\delta)\,C_t$ 是消费的运动方程；与资源约束合起来便给出政策函数。

\ex[吃蛋糕的政策函数]{
在对数效用下，把消费求成存量的函数，并验证它确实解了该问题。设初始总存量为 $H_1 = H$。}
\sol{
猜测消费与存量成比例，$C_t = \gamma\, H_t$，其中 $\gamma$ 是一个待定常数。资源约束随之给出存量的转移，
\[
H_{t+1} = (1-\delta)H_t - C_t = (1-\delta-\gamma)\,H_t,
\]
从而 $C_{t+1} = \gamma\, H_{t+1} = \gamma(1-\delta-\gamma)\,H_t$。施加欧拉方程 $C_{t+1}/C_t = \beta(1-\delta)$，
\[
\frac{\gamma(1-\delta-\gamma)H_t}{\gamma H_t} = 1-\delta-\gamma = \beta(1-\delta)
\quad\Longrightarrow\quad
\gamma = (1-\delta) - \beta(1-\delta) = (1-\beta)(1-\delta).
\]
于是政策函数与由此得到的路径为
\[
C_t = (1-\beta)(1-\delta)\, H_t,
\qquad
H_{t+1} = \beta(1-\delta)\, H_t.
\]
从 $H_1 = H$ 出发，存量以几何方式衰减，$H_t = \big(\beta(1-\delta)\big)^{t-1} H$，所以消费为
\[
C_1 = (1-\beta)(1-\delta)\,H,
\qquad
C_t = (1-\beta)(1-\delta)\,\big(\beta(1-\delta)\big)^{t-1}\, H .
\]
若 $\delta = 0$，这就坍缩成经典的吃蛋糕答案 $C_t = (1-\beta)\,\beta^{\,t-1} H$：每期把剩余的固定比例 $1-\beta$ 吃掉。}

\rmkb[对课堂讲义的一处更正]{讲义把消费比例写成了 $1 - \beta(1-\delta)$，给出 $C_1 = \big(1-\beta(1-\delta)\big)H$。正确的常数是 $(1-\beta)(1-\delta)$。两者仅在特例 $\delta=0$ 时一致，此时都化为 $1-\beta$；当折旧为正时，上面经过验证的答案才是对的——只需直接拿它对照欧拉方程便可核实。}

\subsection{吃蛋糕问题教会了我们什么}

这道习题的价值是概念性的，而非计算性的。有两条教训会延续下去。

第一，辨明谁是状态、谁是控制。这里的存量 $H$ 所扮演的角色，恰恰就是此前各模型里\emph{储蓄}所扮演的角色：它是从过去带到现在的那笔资源。认出这一对应关系，正是让我们把所有这些模型看作同一主题之变奏的关键。

第二，留意约束代表着什么。在前面几节那些受预算约束的问题里，相对价格登场、并支配着各种权衡。在吃蛋糕问题里则根本没有价格——约束是一个纯粹的\emph{资源约束}，只陈述这种物品在物理上有多少可用。这恰恰是当一个单独的决策者（而非市场）来配置资源时、约束所采取的形式。它是通往计划者问题的桥梁，我们现在就转向计划者问题。

\state{概念重于代数}{
动态宏观里反复出现的纪律是：要分清什么是\emph{状态}、什么是\emph{控制}，以及一个约束反映的是\emph{价格}（市场配置）还是纯粹的\emph{资源}（计划者配置）。数学形式是次要的；对状态、控制、约束的经济解读，才是让这些模型可以相互比较的东西。}

\section{社会计划者问题与福利定理}

在竞争性一般均衡中，家庭最大化效用、厂商最大化利润，整个配置由供给与需求通过价格的相互作用所驱动。现在改设想一个仁慈的\emph{社会计划者}，他掌握经济中的全部信息，直接配置资源，根本不用价格。两种配置何时会相同？这个问题的答案是经济学中最深刻的结论之一。

计划者求解
\[
\max_{\set{C_t}_{t=0}^{\infty}}\; \sum_{t=0}^{\infty} \beta^t\, U(C_t)
\]
受制于经济的资源约束
\[
c_t + i_t = y_t,
\]
它说产出被划分为消费与投资——并且关键地请注意，这里没有出现任何价格。产出由资本与劳动、用厂商所用的同一技术生产出来，
\[
y_t = F(k_t, n_t),
\]
而资本通过
\[
k_{t+1} = (1-\delta)\,k_t + i_t,
\]
积累起来，$k_0$ 给定。吃蛋糕问题是 $F \equiv 0$ 的特例，在其中唯一的“资源”就是上一期剩下的存量；计划者问题把它推广开来，让剩下的存量变得\emph{有生产力}——仿佛没吃掉的蛋糕被借了出去，赚来一笔额外的产出 $F(k_t)$。

令 $n_t = 1$，把资源约束与积累方程合起来消去 $i_t$，
\[
c_t + \big(k_{t+1} - (1-\delta)k_t\big) = F(k_t),
\]
于是计划者问题为
\[
\max_{\set{c_t}_{t=0}^{\infty}}\; \sum_{t=0}^{\infty} \beta^t\, U(c_t)
\qquad \text{s.t.} \qquad
c_t + k_{t+1} = F(k_t) + (1-\delta)k_t, \quad k_0 \text{ 给定}.
\]

\subsection{求解计划者问题}

为求清晰，取完全折旧的情形 $\delta = 1$，于是约束读作 $c_t + k_{t+1} = F(k_t)$。拉格朗日函数为
\[
\mathcal{L} = \sum_{t=0}^{\infty} \beta^t\, U(c_t) + \sum_{t=0}^{\infty} \lambda_t\!\pr{F(k_t) - k_{t+1} - c_t},
\]
一阶条件为
\[
\ba
c_t:&\quad \beta^t\, U'(c_t) = \lambda_t, \\
c_{t+1}:&\quad \beta^{t+1}\, U'(c_{t+1}) = \lambda_{t+1}, \\
k_{t+1}:&\quad \lambda_{t+1}\, F'(k_{t+1}) = \lambda_t.
\ea
\]
消去乘子得到计划者的欧拉方程，
\[
\frac{U'(c_t)}{\beta\, U'(c_{t+1})} = F'(k_{t+1}).
\]
这正是吃蛋糕的欧拉方程，只不过把 $1-\delta$ 换成了资本的边际产出 $F'(k_{t+1})$：资源“增长”的速率如今是技术的回报，而不再是一个固定的留存率。

\ex[柯布--道格拉斯的社会计划者]{
在 $F(k_t) = A\, k_t^{\theta}$ 与对数效用 $U(c_t) = \log c_t$ 下，解出资本的运动方程与稳态。}
\sol{
在对数效用下欧拉方程变成
\[
\frac{c_{t+1}}{\beta\, c_t} = F'(k_{t+1}) = A\,\theta\, k_{t+1}^{\theta-1}.
\]
用约束给出的 $c_t = A k_t^{\theta} - k_{t+1}$ 代入，得到一个只含资本的方程，
\[
\frac{A k_{t+1}^{\theta} - k_{t+2}}{\beta\,(A k_t^{\theta} - k_{t+1})} = A\,\theta\, k_{t+1}^{\theta-1}.
\]
这个方程的形式提示我们猜测 $k_{t+1} = B\, k_t^{\theta}$，其中 $B$ 为常数。那么 $c_t = A k_t^{\theta} - B k_t^{\theta} = (A-B)k_t^{\theta}$ 且 $c_{t+1} = (A-B)k_{t+1}^{\theta}$。代入欧拉方程，
\[
\frac{(A-B)k_{t+1}^{\theta}}{\beta (A-B)k_t^{\theta}}
= \frac{B^{\theta}}{\beta}\,k_t^{\theta^2-\theta}
\;\overset{!}{=}\;
A\theta\,(B k_t^{\theta})^{\theta-1}
= A\theta\, B^{\theta-1}\,k_t^{\theta^2-\theta}.
\]
$k_t$ 的幂次相符，证实了猜测；令系数相等，
\[
\frac{B^{\theta}}{\beta} = A\theta\, B^{\theta-1}
\quad\Longrightarrow\quad
B = A\beta\theta,
\]
所以运动方程为
\[
k_{t+1} = A\beta\theta\, k_t^{\theta}.
\]
令 $k_{t+1} = k_t = k_{\mathrm{ss}}$ 并除以 $k_{\mathrm{ss}}^{\theta}$，
\[
k_{\mathrm{ss}}^{\,1-\theta} = A\beta\theta
\quad\Longrightarrow\quad
k_{\mathrm{ss}} = (A\beta\theta)^{\frac{1}{1-\theta}}.
\]
和 OLG 模型一样，生产率水平 $A$ 的增长会把这个目标向上推移，于是稳态沿一条增长路径移动，而非停在原地。}

这个计划者配置与索罗稳态的结构完全相同，但它如今是从跨期\emph{最优化}中导出的，而非从一个假设的储蓄率得来。这正是新古典纲领的全部要旨所在：结论存活了下来，但它们扎根于选择。又因为模型是从明确的偏好与技术搭起来的，它的参数可以被重新诠释以涵盖更多东西——比如，可以通过让生产率随机化来引入经济周期，$y_t = z\, A\, k_t^{\theta}$，其中 $z$ 是随机的（从某个分布中抽取，或在高、低两种状态之间切换）。单一代表性主体的计划者解随后便能推广到整个经济，因为每个主体都彼此相同。

\subsection{两条福利定理}

现在我们可以陈述计划者配置与竞争性均衡之间的关系。在合适的条件下两者重合，而联结它们的桥梁，恰恰是计划者问题里\emph{不}包含的那组价格。计划者的资源约束携带着影子价格——即乘子 $\lambda_t$——而这些影子价格，正是能够把计划者配置去中心化（decentralize）的那些竞争性价格。

\thm{两条福利定理}{
\begin{enumerate}
  \item \textbf{第一福利定理。} 一个竞争性均衡配置是帕累托最优的：在合适的价格下，那个去中心化的结果——每个家庭最大化效用、每个厂商最大化利润——已经无法在不损害某个主体的前提下改善另一个主体了。
  \item \textbf{第二福利定理。} 任何一个帕累托最优配置——尤其是一个仁慈的社会计划者所会选择的那个——都可以被支撑为一个竞争性均衡，只要价格设置得当（并且在必要时，用一次性总额转移来重新分配初始财富）。
\end{enumerate}}

其经济内涵是：市场与计划，尽管组织原则截然相反，在这个无摩擦的世界里却抵达同一个归宿。计划者按资源约束来配置、不理会价格；市场按价格来配置、不理会任何中央指令。然而计划者最优化中隐含的影子价格，恰恰就是那组使家庭与厂商——仅仅出于各自的私利行事——自发选出计划者那套配置的市场价格。正是这一等价关系，许可我们在此后整门课程里去解一个计划者问题（往往要容易得多），并从中读出竞争性均衡。

\rmkb[这是要往哪里去]{新古典模型给了储蓄率一个微观经济学的基础，并借此给了我们一种方法去追问市场结果是否有效率。福利定理在这个基准经济里回答“是”。宏观经济学其余的大部分内容，正是研究当这些定理\emph{失效}时会发生什么——当货币、黏性价格、外部性或缺失的市场在竞争性结果与计划者理想之间打入一道楔子，为政策改善市场腾出空间时。第~\ref{ch:malthus} 章的马尔萨斯经济是第一个这样的背离；货币摩擦与凯恩斯摩擦紧随其后。}
